% BHU PYQ -- Manually transcribed & error-corrected
% Paper: GLB-301: Petrology and Economic Geology
% Exam: B.Sc. (Hons.) Semester III Examination, 2023-24

\documentclass[11pt,a4paper]{article}
\usepackage[a4paper,top=2.5cm,bottom=2.5cm,left=2.8cm,right=2.8cm,headheight=14pt]{geometry}
\usepackage{amsmath,amssymb}
\usepackage[T1]{fontenc}
\usepackage[utf8]{inputenc}
\usepackage{lmodern,microtype,fancyhdr,enumitem,hyperref,lastpage,parskip}

\hypersetup{colorlinks=true,urlcolor=black,linkcolor=black,
  pdftitle={GLB-301 Petrology and Economic Geology -- BHU B.Sc. Sem III 2023-24}}
\pagestyle{fancy}\fancyhf{}
\renewcommand{\headrulewidth}{0.4pt}\renewcommand{\footrulewidth}{0.4pt}
\fancyhead[L]{\small   \textit{Banaras Hindu University}}
\fancyhead[R]{\small   \textit{GLB-301: Petrology \& Economic Geology}}
\fancyfoot[C]{\small Page  \thepage\ of \pageref{LastPage}}


\newlist{parts}{enumerate}{2}
\setlist[parts,1]{label=(\alph*),leftmargin=*,itemsep=5pt,topsep=3pt}

\newcommand{\pts}[1]{\hfill   \textit{[#1]}}


\begin{document}

\begin{center}
  {\large\bfseries BANARAS HINDU UNIVERSITY}\\[2pt]
  {\normalsize B.Sc. (Hons.) --- Semester III Examination, 2023-24}\\[6pt]
  \rule{0.8\linewidth}{0.5pt}\\[6pt]
  {\large\bfseries Paper: GLB-301 --- Petrology and Economic Geology}\\[4pt]
  {\normalsize Time: 3 Hours \hfill Maximum Marks: 70}\\[2pt]
  {\small Ref. No.: 067528}
\end{center}
\rule{\linewidth}{0.4pt}
\smallskip



\noindent   \textit{Instructions: Answer   \textbf{five} questions in all, selecting at least   \textbf{two} each from Sections A and B, and Section-C which is compulsory.}
\bigskip

\bigskip


\noindent {\large\bfseries SECTION --- A}
\rule{\linewidth}{0.4pt}\smallskip




\noindent   \textbf{Question 1.} Write a brief note on the phase rule application to the $ \text{H}_2 \text{O}$ system with a suitable diagram. \pts{14}

\medskip


\noindent   \textbf{Question 2.} Discuss the classification of sedimentary rocks. \pts{14}

\medskip


\noindent   \textbf{Question 3.} Briefly describe metamorphic textures and structures. \pts{14}

\medskip


\noindent   \textbf{Question 4.} Write short notes on   \textbf{any four} of the following: \pts{$4  \times 3\frac{1}{2} = 14$}

\begin{parts}
  \item Bowen's reaction series
  \item Batholith
  \item Conglomerate
  \item Shale
  \item Progressive metamorphism
  \item Hornfels
\end{parts}

\bigskip


\noindent {\large\bfseries SECTION --- B}
\rule{\linewidth}{0.4pt}\smallskip




\noindent   \textbf{Question 5.} Write important chemical formula, distribution in India and uses of   \textbf{any four} of the following economic minerals: \pts{$4  \times 3\frac{1}{2} = 14$}

\begin{parts}
  \item Halite
  \item Quartz
  \item Biotite
  \item Pyrite
  \item Malachite
  \item Monazite
\end{parts}

\medskip


\noindent   \textbf{Question 6.} Write in brief   \textbf{any two} of the following: \pts{$2  \times 7 = 14$}

\begin{parts}
  \item What is petroleum? Discuss the origin and economic uses of petroleum.
  \item What are fertilizer minerals? Discuss the economic uses of fertilizer minerals and their distribution in India.
  \item What are galena and sphalerite? Describe the economic uses and distribution of lead-zinc ore deposits in India.
\end{parts}

\medskip


\noindent   \textbf{Question 7.} Briefly describe the physical properties of manganese ore minerals, economic uses, and distribution of mineral deposits in the map of India. \pts{14}

\bigskip


\noindent {\large\bfseries SECTION --- C}
\rule{\linewidth}{0.4pt}\smallskip




\noindent   \textbf{Question 8.} Answer the following:

\begin{parts}
  \item What is the mineralogical composition of peridotite? \pts{1}
  \item Differentiate between breccia and conglomerate. \pts{2}
  \item Define the matrix of sedimentary rocks. \pts{2}
  \item What is slate? \pts{1}
  \item Write the name of a low-grade metamorphic rock. \pts{1}
  \item Write the names of two coal deposits in India. \pts{2}
  \item Write the name of an ore mineral of beryllium. \pts{1}
  \item How many sets of cleavage planes are found in kyanite? \pts{1}
  \item Which metal is extracted from magnesite? \pts{1}
  \item What are the economic uses of chalcopyrite? \pts{2}
\end{parts}

\vfill\rule{\linewidth}{0.4pt}

\begin{center}\small
\href{https://bhu-exam.vercel.app/}{Click here to visit our website}\\[4pt]
\href{https://me-aryan.vercel.app/}{Click here to visit our contributor}\\[4pt]
\href{https://quashan-me.vercel.app/}{Click here to visit our contributor}
\end{center}
\end{document}
